\subsection{Glossary}
\label{subsection_glossary}

\textbf{Moment and torque} (source Wikipedia):
\begin{itemize}
    \item A moment is a mathematical expression involving the product of a distance and a
    physical quantity such as a force or electric charge. Moments are usually defined with
    respect to a fixed reference point and refer to physical quantities located some
    distance from the reference point. The moment of force, often called torque, is the
    product of a force on an object and the distance from the reference point to the
    object.

    \item In physics and mechanics, torque is the rotational correspondent of linear
    force. It is referred to as the moment of force, or simply the moment. Just as a
    linear force is a push or a pull applied to a body, a torque can be thought of as a
    twist applied to an object with respect to a chosen point.

    \item Torque is referred to using different vocabulary depending on geographical
    location and field of study. Torque and moment are often used interchangeably for the
    rotational effect of a force, but in mechanics, the main difference is application:
    torque refers to a force causing rotation in a dynamic system, while moment is often
    used for the turning effect in a static, or non-moving, system. Both are calculated as
    the force multiplied by the perpendicular distance from the pivot point.

    \item When modelling physics in $PGA$ forces and moments, describing the translational
    and rotational physical aspects, do not need to be considered separately, but can be
    modelled within one object containing all information. For example, by modelling the
    line of action $F$ of a force $f$ acting at a point $P$ as $F = P \wedge f$ as in
    equation~(\ref{eq:line_of_action}) both effects are covered, the force $f$ itself and
    the torque generated at any point are covered by the mathematical object $F$. The
    torque generated at the origin $M_{O} = bulk(F)$ can be as easily recovered from $F$
    as the resulting free moment at any other point.
\end{itemize}

\textbf{Momentum}
\begin{itemize}
    \item Momentum (German: Impuls, Bewegungsgröße) characterizes the movement of an
    object. It has a direction (the direction of the speed vector) and a magnitude (the
    product of the mass and the magnitude of the speed vector). A heavier object moving at
    the same speed has more momentum compared to an object of less mass. Momentum can
    change in direction and/or magnitude.
\end{itemize}

\textbf{Pose}
\begin{itemize}
    \item A pose is the rigid placement of one frame relative to another: an origin and an
    orientation (in $2D$ a position and a scalar angle $\varphi$; in $3D$ a position and
    an axis-angle vector).

    \item It is a point of the rigid-motion group, carried by a unit motor $M$; the
    origin/orientation form and the motor interconvert exactly
    (see Sec.~\ref{subsubsection:pose_frame_tree}).

    \item A pose is a finite placement (where a frame is), as opposed to a twist (how it
    is moving).
\end{itemize}

\textbf{Motor and rotor}
\begin{itemize}
    \item In $PGA$ a motor $M$ performs a rigid motion by the sandwich product $X' = M
    \veedot X \veedot \rrev{M}$, applied identically to every object: points, lines,
    twists, wrenches.

    \item It is the exponential of a twist, $M = \exp_\veedot(\frac{1}{2} B)$ (see
    eqn.~(\ref{eq:motor_pga_body_frame})); motors multiply (compose) where their
    generators add.

    \item A unit motor in $3D$ $PGA$ encodes a screw motion (rotation about and
    translation along an axis), so one motor covers rotation, translation and their
    combination.

    \item The $EGA$ counterpart, built from the geometric product $\wedgedot$ instead of
    the regressive $\veedot$ and acting only through the origin, is the rotor (rotations
    only).
\end{itemize}

\textbf{Twist}
\begin{itemize}
    \item A twist $B$ is the instantaneous motion of a frame: the Lie-algebra generator
    that builds its motor through $M = \exp_\veedot(\frac{1}{2} B)$ (see
    eqn.~(\ref{eq:motor_pga_body_frame})). Its grade is dimension-dependent: a vector in
    $2D$ $PGA$, a bivector in $3D$ $PGA$.

    \item The rate $\Omega$ is the time derivative of the generator, $\Omega = \dot{B} =
    dB/dt$. For a constant rate the generator grows linearly, $B(t) = \Omega\,t$, so the
    motor is $M(t) = \exp_\veedot(\frac{1}{2}\,\Omega\,t)$; conversely the rate is recovered
    from a moving motor by $\Omega = 2\,\dot{M} \veedot \rrev{M}$ (see
    eqn.~(\ref{eq:motor_rate_body_frame_pga})).

    \item Twists add (they live in a vector space) where motors multiply.

    \item A twist is a velocity screw; the point velocity it induces is $\dot{X} =
    \mathrm{rcmt}(\Omega, X)$.
\end{itemize}

\textbf{Pose vs.\ twist}
\begin{itemize}
    \item Pose is the group element (where a frame is, the motor $M$); twist is the
    tangent/rate (how it moves, the generator $B$ with rate $\Omega = \dot{B}$).

    \item They interconvert by integration and differentiation: $M(t) = M_0 \veedot
    \exp_\veedot(\frac{1}{2} \int \Omega\,dt)$ and conversely $\Omega = 2\,\dot{M} \veedot
    \rrev{M}$.

    \item Pose lives in the Lie group, twist in the Lie algebra (see
    Sec.~\ref{subsubsection_lie_algebra_vs_lie_group}); cf.\ the entry on infinitesimal
    vs.\ finite movements.
\end{itemize}

\textbf{Screw (Chasles' theorem)}
\begin{itemize}
    \item By Chasles' theorem every rigid motion is a screw: a rotation about an axis
    combined with a translation along it.

    \item In $PGA$ the screw axis is a line (a bivector in $3D$); the central quantities
    are all screws on such an axis: the twist (a velocity screw), the wrench (a force
    screw), and the motor (the finite screw motion $\exp_\veedot$ of a twist).

    \item A pure rotation and a pure translation are the degenerate screws of zero and
    infinite pitch.
\end{itemize}

\textbf{Adjoint}
\begin{itemize}
    \item The adjoint $\mathrm{Ad}_M$ re-expresses a twist in another frame.

    \item In $PGA$ it is the same sandwich that moves any object, $\mathrm{Ad}_M(B) = M
    \veedot B \veedot \rrev{M}$ (see eqn.~(\ref{eq:adjoint_is_sandwich})).

    \item Point action and twist transport are thus one operation; no separate $6\times6$
    adjoint matrix is needed (see Sec.~\ref{subsubsection:transport_along_chain}).
\end{itemize}

\textbf{Commutator (Lie bracket)}
\begin{itemize}
    \item The Lie bracket (commutator) of two twists measures how far they fail to
    commute.

    \item In $PGA$ it is a native product, $[A, B] = \mathrm{rcmt}(A, B)$: the
    $\mathfrak{se}(2)/\mathfrak{se}(3)$ twist bracket.

    \item It supplies the Coriolis/centrifugal coupling in the Newton-Euler recursion (see
    eqn.~(\ref{eq:newton_euler_recursion})) and the gyroscopic term $\mathrm{rcmt}(\Omega,
    I[\Omega])$ in the body-frame Euler equation.
\end{itemize}

\textbf{Lie group vs.\ Lie algebra}
\begin{itemize}
    \item The motors form a Lie group: the manifold of finite rigid motions, composed by
    $\veedot$.

    \item The twists form its Lie algebra: the tangent space at the identity, a vector
    space closed under the bracket $[A,B] = \mathrm{rcmt}(A,B)$.

    \item The two are linked by $\exp_\veedot$ (algebra to group) and its logarithm (group
    to algebra); twists add and commute to first order, motors multiply and do not (see
    Sec.~\ref{subsubsection_lie_algebra_vs_lie_group}).
\end{itemize}

\textbf{Infinitesimal vs.\ finite movements}
\begin{itemize}
    \item An infinitesimal movement is a twist (a Lie-algebra generator); a finite
    movement is the motor it integrates to (a Lie-group element).

    \item Simultaneous infinitesimal motions superpose (twists add), but the finite motors
    they generate do not commute and must be multiplied in order.

    \item This is why a parent-to-child pose is built by composing a rotation and a
    translation motor, not by exponentiating their summed generators (see
    Sec.~\ref{subsubsection_lie_algebra_vs_lie_group}).
\end{itemize}

\textbf{Wrench}
\begin{itemize}
    \item A wrench is the force/torque screw dual to a twist: the bivector line of action
    carrying a force together with the moment it generates (the ``forque'' of
    Sec.~\ref{subsubsection:modelling_force_and_torque}).

    \item The momentum line of Sec.~\ref{subsubsection:modelling_momentum} is likewise the
    wrench dual to the velocity twist.

    \item A twist $\xi$ pairs with a wrench $P$ through the reciprocal (spatial) product
    $\langle \xi, P \rangle = -\,\mathrm{rwdg}(\xi, P)$, giving power / generalised force;
    with $P = I[\xi]$ it gives twice the kinetic energy (see
    eqn.~(\ref{eq:spatial_pairing})).
\end{itemize}

\textbf{Inertia map}
\begin{itemize}
    \item The inertia map $I[\cdot]$ is the linear operator sending a body's velocity
    twist $\Omega$ to its momentum wrench, $P = I[\Omega]$ (see
    Sec.~\ref{subsubsection:structure_inertia_maps}).

    \item It is a $3\times3$ map (vector to bivector) in $2D$ and a $6\times6$ map
    (bivector to bivector) in $3D$, encoding mass, scalar moment of inertia and the full
    inertia tensor in one object; the mass contribution is $-m\,\rcmpl{\Omega}$, the
    complement placing it in the dual block.

    \item It is additive over the parts of a body and is inverted to set up the Euler
    equation $\dot{\Omega} = I^{-1}[\,W - \mathrm{rcmt}(\Omega, I[\Omega])\,]$.
\end{itemize}

\newpage
